% Môn: Vật lý
% Lớp: 12
% Chương: Chương II. Khí lí tưởng
% Bài: L12C2 Bài 13. Bài tập về khí lí tưởng
% Số câu: 132
% App đọc file này trực tiếp — sửa rồi Commit trên GitHub, bấm Đồng bộ GitHub .tex

% L12C2 Bài 13. Bài tập về khí lí tưởng

% ===== Câu 1 =====
% ID: L12C2B13-01-DS
% Mức: TH
\dangbt{Bài toán phương trình trạng thái và Clapeyron}
\begin{ex}
% ID: L12C2B13-01-DS
% Mức: TH
Xét dạng “Bài toán phương trình trạng thái và Clapeyron”: a) Công thức đặc trưng là $pV=nRT$ và $p_1V_1/T_1=p_2V_2/T_2$. b) Có thể bỏ qua điều kiện khí lí tưởng, các đại lượng dùng theo SI. c) trạng thái khí được xác định bởi p, $V$, $T$ và lượng khí. d) Có thể không đổi lít sang m³ hoặc Celsius sang kelvin.
\choiceTF

\loigiai{
a) Đúng: phù hợp với công thức và điều kiện đã nêu. b) Sai: phát biểu không thỏa công thức hoặc điều kiện áp dụng. c) Đúng: phù hợp với công thức và điều kiện đã nêu. d) Sai: phát biểu không thỏa công thức hoặc điều kiện áp dụng.
}
\end{ex}

% ===== Câu 2 =====
% ID: L12C2B13-01-TLN
% Mức: TH
\begin{ex}
% ID: L12C2B13-01-TLN
% Mức: TH
Có 2 mol khí lí tưởng ở 300 $K$ chiếm 0,050 m³. Lấy $R$ = $8,31\,\text{J}$ /(mol·K). Tính áp suất gần đúng theo kPa.
\shortans{99,7}
\loigiai{
$p=nRT/V=2\cdot8,31\cdot300/0,050=99720$ Pa ≈ $99,7\,\text{kPa}$.
}
\end{ex}

% ===== Câu 3 =====
% ID: L12C2B13-01-TN
% Mức: NB
\begin{ex}
% ID: L12C2B13-01-TN
% Mức: NB
Công thức cốt lõi phù hợp nhất cho dạng “Bài toán phương trình trạng thái và Clapeyron” là
\choice
{\True $pV=nRT$ và $p_1V_1/T_1=p_2V_2/T_2$}
{$Q=mc\Delta T$}
{$A=Fs$}
{$U=RI$}
\loigiai{
Đáp án A. Dùng $pV=nRT$ và $p_1V_1/T_1=p_2V_2/T_2$ trong điều kiện khí lí tưởng, các đại lượng dùng theo SI.
}
\end{ex}

% ===== Câu 4 =====
% ID: L12C2B13-02-DS
% Mức: TH
\begin{ex}
% ID: L12C2B13-02-DS
% Mức: TH
Xét dạng “Bài toán phương trình trạng thái và Clapeyron”: a) Cần dùng đơn vị phù hợp: Pa; m³; mol; K. b) Nhiệt độ tuyệt đối có thể mang giá trị âm. c) Phải xác định đại lượng không đổi trước khi tính. d) Không cần kiểm tra thứ nguyên.
\choiceTF

\loigiai{
a) Đúng: phù hợp với công thức và điều kiện đã nêu. b) Sai: phát biểu không thỏa công thức hoặc điều kiện áp dụng. c) Đúng: phù hợp với công thức và điều kiện đã nêu. d) Sai: phát biểu không thỏa công thức hoặc điều kiện áp dụng.
}
\end{ex}

% ===== Câu 5 =====
% ID: L12C2B13-02-TLN
% Mức: TH
\begin{ex}
% ID: L12C2B13-02-TLN
% Mức: TH
Đổi 27 ° $C$ sang nhiệt độ tuyệt đối theo kelvin.
\shortans{300}
\loigiai{
$T=27+273=300$ K.
}
\end{ex}

% ===== Câu 6 =====
% ID: L12C2B13-02-TN
% Mức: NB
\begin{ex}
% ID: L12C2B13-02-TN
% Mức: NB
Điều kiện quan trọng khi áp dụng dạng “Bài toán phương trình trạng thái và Clapeyron” là
\choice
{khối lượng vật luôn bằng $1\,\text{kg}$}
{\True khí lí tưởng, các đại lượng dùng theo SI}
{nhiệt độ luôn bằng 0 ° $C$}
{áp suất luôn bằng áp suất khí quyển}
\loigiai{
Đáp án B. Điều kiện áp dụng là khí lí tưởng, các đại lượng dùng theo SI.
}
\end{ex}

% ===== Câu 7 =====
% ID: L12C2B13-03-DS
% Mức: VD
\begin{ex}
% ID: L12C2B13-03-DS
% Mức: VD
Xét dạng “Bài toán phương trình trạng thái và Clapeyron”: a) Bài mẫu cho kết quả 99,7. b) Kết quả bài mẫu gấp đôi 99,7. c) Lời giải phải nêu công thức trước khi thay số. d) Mọi dữ kiện số đều không cần đổi đơn vị.
\choiceTF

\loigiai{
a) Đúng: phù hợp với công thức và điều kiện đã nêu. b) Sai: phát biểu không thỏa công thức hoặc điều kiện áp dụng. c) Đúng: phù hợp với công thức và điều kiện đã nêu. d) Sai: phát biểu không thỏa công thức hoặc điều kiện áp dụng.
}
\end{ex}

% ===== Câu 8 =====
% ID: L12C2B13-03-TLN
% Mức: VD
\begin{ex}
% ID: L12C2B13-03-TLN
% Mức: VD
Một đại lượng tăng từ 100 lên 150. Tính tỉ số giá trị sau so với trước.
\shortans{1,5}
\loigiai{
Tỉ số bằng $150/100=1,5$.
}
\end{ex}

% ===== Câu 9 =====
% ID: L12C2B13-03-TN
% Mức: TH
\begin{ex}
% ID: L12C2B13-03-TN
% Mức: TH
Nhận xét đúng về quan hệ các đại lượng là
\choice
{các đại lượng không phụ thuộc nhau}
{mọi đại lượng đều không đổi}
{\True trạng thái khí được xác định bởi p, $V$, $T$ và lượng khí}
{chỉ phụ thuộc thời gian}
\loigiai{
Đáp án C. trạng thái khí được xác định bởi p, $V$, $T$ và lượng khí.
}
\end{ex}

% ===== Câu 10 =====
% ID: L12C2B13-04-DS
% Mức: VD
\begin{ex}
% ID: L12C2B13-04-DS
% Mức: VD
Xét dạng “Bài toán phương trình trạng thái và Clapeyron”: a) Quan hệ cần dùng là trạng thái khí được xác định bởi p, $V$, $T$ và lượng khí. b) Có thể suy ra quan hệ mà không xét điều kiện quá trình. c) Kết quả phải phù hợp chiều biến thiên vật lí. d) Áp suất dư luôn thay được áp suất tuyệt đối.
\choiceTF

\loigiai{
a) Đúng: phù hợp với công thức và điều kiện đã nêu. b) Sai: phát biểu không thỏa công thức hoặc điều kiện áp dụng. c) Đúng: phù hợp với công thức và điều kiện đã nêu. d) Sai: phát biểu không thỏa công thức hoặc điều kiện áp dụng.
}
\end{ex}

% ===== Câu 11 =====
% ID: L12C2B13-04-TLN
% Mức: VD
\begin{ex}
% ID: L12C2B13-04-TLN
% Mức: VD
Một thể tích 2,5 $L$ bằng bao nhiêu m³? Viết kết quả theo dạng $a\cdot10^{-3}$.
\shortans{2,5}
\loigiai{
$2,5L=2,5\cdot10^{-3}$ m³.
}
\end{ex}

% ===== Câu 12 =====
% ID: L12C2B13-04-TN
% Mức: TH
\begin{ex}
% ID: L12C2B13-04-TN
% Mức: TH
Sai lầm cần tránh khi giải dạng này là
\choice
{ghi đầy đủ dữ kiện}
{đổi đơn vị đúng}
{kiểm tra điều kiện}
{\True không đổi lít sang m³ hoặc Celsius sang kelvin}
\loigiai{
Đáp án D. Cần tránh: không đổi lít sang m³ hoặc Celsius sang kelvin.
}
\end{ex}

% ===== Câu 13 =====
% ID: L12C2B13-05-DS
% Mức: VDC
\begin{ex}
% ID: L12C2B13-05-DS
% Mức: VDC
Xét dạng “Bài toán phương trình trạng thái và Clapeyron”: a) Dạng này có thể yêu cầu phối hợp đổi đơn vị và lập tỉ số. b) Kelvin được đổi sang Celsius bằng cách cộng 273. c) Cần ghi rõ đại lượng cần tìm. d) Công thức nào có đủ kí hiệu cũng dùng được.
\choiceTF

\loigiai{
a) Đúng: phù hợp với công thức và điều kiện đã nêu. b) Sai: phát biểu không thỏa công thức hoặc điều kiện áp dụng. c) Đúng: phù hợp với công thức và điều kiện đã nêu. d) Sai: phát biểu không thỏa công thức hoặc điều kiện áp dụng.
}
\end{ex}

% ===== Câu 14 =====
% ID: L12C2B13-05-TLN
% Mức: VD
\begin{ex}
% ID: L12C2B13-05-TLN
% Mức: VD
Áp suất $120\,\text{kPa}$ bằng bao nhiêu Pa, lấy hệ số của $10^3$ làm đáp án?
\shortans{120}
\loigiai{
$120\,\text{kPa}=120\cdot10^3$ Pa.
}
\end{ex}

% ===== Câu 15 =====
% ID: L12C2B13-05-TN
% Mức: TH
\begin{ex}
% ID: L12C2B13-05-TN
% Mức: TH
Đơn vị thường dùng sau khi chuẩn hóa theo SI là
\choice
{\True Pa; m³; mol; $K$}
{kg·m}
{W/s}
{C/mol}
\loigiai{
Đáp án A. Các đơn vị phù hợp: Pa; m³; mol; K.
}
\end{ex}

% ===== Câu 16 =====
% ID: L12C2B13-06-DS
% Mức: TH
\dangbt{Bài toán quá trình đẳng nhiệt}
\begin{ex}
% ID: L12C2B13-06-DS
% Mức: TH
Xét dạng “Bài toán quá trình đẳng nhiệt”: a) Công thức đặc trưng là $p_1V_1=p_2V_2$. b) Có thể bỏ qua điều kiện lượng khí và nhiệt độ không đổi. c) p tỉ lệ nghịch với V. d) Có thể kết luận p tỉ lệ thuận với V.
\choiceTF

\loigiai{
a) Đúng: phù hợp với công thức và điều kiện đã nêu. b) Sai: phát biểu không thỏa công thức hoặc điều kiện áp dụng. c) Đúng: phù hợp với công thức và điều kiện đã nêu. d) Sai: phát biểu không thỏa công thức hoặc điều kiện áp dụng.
}
\end{ex}

% ===== Câu 17 =====
% ID: L12C2B13-06-TLN
% Mức: VDC
\dangbt{Bài toán phương trình trạng thái và Clapeyron}
\begin{ex}
% ID: L12C2B13-06-TLN
% Mức: VDC
Theo quan hệ “trạng thái khí được xác định bởi p, $V$, $T$ và lượng khí”, nếu đại lượng phụ thuộc tỉ lệ thuận tăng gấp 4 thì hệ số tăng là bao nhiêu?
\shortans{4}
\loigiai{
Với quan hệ tỉ lệ thuận, hệ số biến đổi được giữ nguyên: 4.
}
\end{ex}

% ===== Câu 18 =====
% ID: L12C2B13-06-TN
% Mức: VD
\begin{ex}
% ID: L12C2B13-06-TN
% Mức: VD
Bước đầu tiên hợp lí khi giải bài là
\choice
{thay số ngay}
{\True xác định quá trình và đại lượng không đổi}
{bỏ qua đơn vị}
{đổi áp suất thành nhiệt độ}
\loigiai{
Đáp án B. Trước hết phải nhận diện quá trình, liệt kê dữ kiện và đại lượng không đổi.
}
\end{ex}

% ===== Câu 19 =====
% ID: L12C2B13-07-DS
% Mức: TH
\dangbt{Bài toán quá trình đẳng nhiệt}
\begin{ex}
% ID: L12C2B13-07-DS
% Mức: TH
Xét dạng “Bài toán quá trình đẳng nhiệt”: a) Cần dùng đơn vị phù hợp: Pa; m³. b) Nhiệt độ tuyệt đối có thể mang giá trị âm. c) Phải xác định đại lượng không đổi trước khi tính. d) Không cần kiểm tra thứ nguyên.
\choiceTF

\loigiai{
a) Đúng: phù hợp với công thức và điều kiện đã nêu. b) Sai: phát biểu không thỏa công thức hoặc điều kiện áp dụng. c) Đúng: phù hợp với công thức và điều kiện đã nêu. d) Sai: phát biểu không thỏa công thức hoặc điều kiện áp dụng.
}
\end{ex}

% ===== Câu 20 =====
% ID: L12C2B13-07-TLN
% Mức: TH
\begin{ex}
% ID: L12C2B13-07-TLN
% Mức: TH
Khí ở $100\,\text{kPa}$ có thể tích 6 $L$, bị nén đẳng nhiệt còn 4 L. Tính áp suất sau theo kPa.
\shortans{150}
\loigiai{
$p_2=p_1V_1/V_2=100\cdot6/4=150$ kPa.
}
\end{ex}

% ===== Câu 21 =====
% ID: L12C2B13-07-TN
% Mức: VD
\dangbt{Bài toán phương trình trạng thái và Clapeyron}
\begin{ex}
% ID: L12C2B13-07-TN
% Mức: VD
Với dữ kiện của bài mẫu: Có 2 mol khí lí tưởng ở 300 $K$ chiếm 0,050 m³. Lấy $R$ = $8,31\,\text{J}$ /(mol·K). Tính áp suất gần đúng theo kPa.
\choice
{\True Kết quả 99,7}
{Kết quả 199.4}
{Không đủ dữ kiện}
{Kết quả bằng 0}
\loigiai{
Đáp án A. $p=nRT/V=2\cdot8,31\cdot300/0,050=99720$ Pa ≈ $99,7\,\text{kPa}$.
}
\end{ex}

% ===== Câu 22 =====
% ID: L12C2B13-08-DS
% Mức: VD
\dangbt{Bài toán quá trình đẳng nhiệt}
\begin{ex}
% ID: L12C2B13-08-DS
% Mức: VD
Xét dạng “Bài toán quá trình đẳng nhiệt”: a) Bài mẫu cho kết quả 150. b) Kết quả bài mẫu gấp đôi 150. c) Lời giải phải nêu công thức trước khi thay số. d) Mọi dữ kiện số đều không cần đổi đơn vị.
\choiceTF

\loigiai{
a) Đúng: phù hợp với công thức và điều kiện đã nêu. b) Sai: phát biểu không thỏa công thức hoặc điều kiện áp dụng. c) Đúng: phù hợp với công thức và điều kiện đã nêu. d) Sai: phát biểu không thỏa công thức hoặc điều kiện áp dụng.
}
\end{ex}

% ===== Câu 23 =====
% ID: L12C2B13-08-TLN
% Mức: TH
\begin{ex}
% ID: L12C2B13-08-TLN
% Mức: TH
Đổi 27 ° $C$ sang nhiệt độ tuyệt đối theo kelvin.
\shortans{300}
\loigiai{
$T=27+273=300$ K.
}
\end{ex}

% ===== Câu 24 =====
% ID: L12C2B13-08-TN
% Mức: VD
\dangbt{Bài toán phương trình trạng thái và Clapeyron}
\begin{ex}
% ID: L12C2B13-08-TN
% Mức: VD
Khi kiểm tra kết quả, tiêu chí vật lí quan trọng là
\choice
{chỉ cần số dương}
{không cần đơn vị}
{\True phù hợp chiều biến thiên theo định luật}
{luôn phải lớn hơn giá trị đầu}
\loigiai{
Đáp án C. Kết quả phải phù hợp với nhận xét: trạng thái khí được xác định bởi p, $V$, $T$ và lượng khí.
}
\end{ex}

% ===== Câu 25 =====
% ID: L12C2B13-09-DS
% Mức: VD
\dangbt{Bài toán quá trình đẳng nhiệt}
\begin{ex}
% ID: L12C2B13-09-DS
% Mức: VD
Xét dạng “Bài toán quá trình đẳng nhiệt”: a) Quan hệ cần dùng là p tỉ lệ nghịch với V. b) Có thể suy ra quan hệ mà không xét điều kiện quá trình. c) Kết quả phải phù hợp chiều biến thiên vật lí. d) Áp suất dư luôn thay được áp suất tuyệt đối.
\choiceTF

\loigiai{
a) Đúng: phù hợp với công thức và điều kiện đã nêu. b) Sai: phát biểu không thỏa công thức hoặc điều kiện áp dụng. c) Đúng: phù hợp với công thức và điều kiện đã nêu. d) Sai: phát biểu không thỏa công thức hoặc điều kiện áp dụng.
}
\end{ex}

% ===== Câu 26 =====
% ID: L12C2B13-09-TLN
% Mức: VD
\begin{ex}
% ID: L12C2B13-09-TLN
% Mức: VD
Một đại lượng tăng từ 100 lên 150. Tính tỉ số giá trị sau so với trước.
\shortans{1,5}
\loigiai{
Tỉ số bằng $150/100=1,5$.
}
\end{ex}

% ===== Câu 27 =====
% ID: L12C2B13-09-TN
% Mức: VDC
\dangbt{Bài toán phương trình trạng thái và Clapeyron}
\begin{ex}
% ID: L12C2B13-09-TN
% Mức: VDC
Nếu giữ đúng các điều kiện của định luật, biểu thức nên dùng là
\choice
{$s=vt$}
{\True $pV=nRT$ và $p_1V_1/T_1=p_2V_2/T_2$}
{$P=A/t$}
{$F=ma$}
\loigiai{
Đáp án B. Biểu thức đặc trưng là $pV=nRT$ và $p_1V_1/T_1=p_2V_2/T_2$.
}
\end{ex}

% ===== Câu 28 =====
% ID: L12C2B13-10-DS
% Mức: VDC
\dangbt{Bài toán quá trình đẳng nhiệt}
\begin{ex}
% ID: L12C2B13-10-DS
% Mức: VDC
Xét dạng “Bài toán quá trình đẳng nhiệt”: a) Dạng này có thể yêu cầu phối hợp đổi đơn vị và lập tỉ số. b) Kelvin được đổi sang Celsius bằng cách cộng 273. c) Cần ghi rõ đại lượng cần tìm. d) Công thức nào có đủ kí hiệu cũng dùng được.
\choiceTF

\loigiai{
a) Đúng: phù hợp với công thức và điều kiện đã nêu. b) Sai: phát biểu không thỏa công thức hoặc điều kiện áp dụng. c) Đúng: phù hợp với công thức và điều kiện đã nêu. d) Sai: phát biểu không thỏa công thức hoặc điều kiện áp dụng.
}
\end{ex}

% ===== Câu 29 =====
% ID: L12C2B13-10-TLN
% Mức: VD
\begin{ex}
% ID: L12C2B13-10-TLN
% Mức: VD
Một thể tích 2,5 $L$ bằng bao nhiêu m³? Viết kết quả theo dạng $a\cdot10^{-3}$.
\shortans{2,5}
\loigiai{
$2,5L=2,5\cdot10^{-3}$ m³.
}
\end{ex}

% ===== Câu 30 =====
% ID: L12C2B13-10-TN
% Mức: VDC
\dangbt{Bài toán phương trình trạng thái và Clapeyron}
\begin{ex}
% ID: L12C2B13-10-TN
% Mức: VDC
Phát biểu nào phản ánh đúng bản chất bài toán?
\choice
{Có thể bỏ qua nhiệt độ tuyệt đối}
{Không cần xác định lượng khí}
{Mọi quá trình đều đẳng nhiệt}
{\True Phải dùng $pV=nRT$ và $p_1V_1/T_1=p_2V_2/T_2$ và kiểm tra khí lí tưởng, các đại lượng dùng theo SI}
\loigiai{
Đáp án D. Phải dùng $pV=nRT$ và $p_1V_1/T_1=p_2V_2/T_2$, đồng thời kiểm tra khí lí tưởng, các đại lượng dùng theo SI.
}
\end{ex}

% ===== Câu 31 =====
% ID: L12C2B13-11-DS
% Mức: TH
\dangbt{Bài toán quá trình đẳng tích}
\begin{ex}
% ID: L12C2B13-11-DS
% Mức: TH
Xét dạng “Bài toán quá trình đẳng tích”: a) Công thức đặc trưng là $p_1/T_1=p_2/T_2$. b) Có thể bỏ qua điều kiện lượng khí và thể tích không đổi. c) p tỉ lệ thuận với $T$ tuyệt đối. d) Có thể bỏ qua việc đổi nhiệt độ sang kelvin.
\choiceTF

\loigiai{
a) Đúng: phù hợp với công thức và điều kiện đã nêu. b) Sai: phát biểu không thỏa công thức hoặc điều kiện áp dụng. c) Đúng: phù hợp với công thức và điều kiện đã nêu. d) Sai: phát biểu không thỏa công thức hoặc điều kiện áp dụng.
}
\end{ex}

% ===== Câu 32 =====
% ID: L12C2B13-11-TLN
% Mức: VD
\dangbt{Bài toán quá trình đẳng nhiệt}
\begin{ex}
% ID: L12C2B13-11-TLN
% Mức: VD
Áp suất $120\,\text{kPa}$ bằng bao nhiêu Pa, lấy hệ số của $10^3$ làm đáp án?
\shortans{120}
\loigiai{
$120\,\text{kPa}=120\cdot10^3$ Pa.
}
\end{ex}

% ===== Câu 33 =====
% ID: L12C2B13-11-TN
% Mức: NB
\begin{ex}
% ID: L12C2B13-11-TN
% Mức: NB
Công thức cốt lõi phù hợp nhất cho dạng “Bài toán quá trình đẳng nhiệt” là
\choice
{\True $p_1V_1=p_2V_2$}
{$Q=mc\Delta T$}
{$A=Fs$}
{$U=RI$}
\loigiai{
Đáp án A. Dùng $p_1V_1=p_2V_2$ trong điều kiện lượng khí và nhiệt độ không đổi.
}
\end{ex}

% ===== Câu 34 =====
% ID: L12C2B13-12-DS
% Mức: TH
\dangbt{Bài toán quá trình đẳng tích}
\begin{ex}
% ID: L12C2B13-12-DS
% Mức: TH
Xét dạng “Bài toán quá trình đẳng tích”: a) Cần dùng đơn vị phù hợp: $K$; Pa. b) Nhiệt độ tuyệt đối có thể mang giá trị âm. c) Phải xác định đại lượng không đổi trước khi tính. d) Không cần kiểm tra thứ nguyên.
\choiceTF

\loigiai{
a) Đúng: phù hợp với công thức và điều kiện đã nêu. b) Sai: phát biểu không thỏa công thức hoặc điều kiện áp dụng. c) Đúng: phù hợp với công thức và điều kiện đã nêu. d) Sai: phát biểu không thỏa công thức hoặc điều kiện áp dụng.
}
\end{ex}

% ===== Câu 35 =====
% ID: L12C2B13-12-TLN
% Mức: VDC
\dangbt{Bài toán quá trình đẳng nhiệt}
\begin{ex}
% ID: L12C2B13-12-TLN
% Mức: VDC
Theo quan hệ “p tỉ lệ nghịch với V”, nếu đại lượng phụ thuộc tỉ lệ thuận tăng gấp 4 thì hệ số tăng là bao nhiêu?
\shortans{4}
\loigiai{
Với quan hệ tỉ lệ thuận, hệ số biến đổi được giữ nguyên: 4.
}
\end{ex}

% ===== Câu 36 =====
% ID: L12C2B13-12-TN
% Mức: NB
\begin{ex}
% ID: L12C2B13-12-TN
% Mức: NB
Điều kiện quan trọng khi áp dụng dạng “Bài toán quá trình đẳng nhiệt” là
\choice
{khối lượng vật luôn bằng $1\,\text{kg}$}
{\True lượng khí và nhiệt độ không đổi}
{nhiệt độ luôn bằng 0 ° $C$}
{áp suất luôn bằng áp suất khí quyển}
\loigiai{
Đáp án B. Điều kiện áp dụng là lượng khí và nhiệt độ không đổi.
}
\end{ex}

% ===== Câu 37 =====
% ID: L12C2B13-13-DS
% Mức: VD
\dangbt{Bài toán quá trình đẳng tích}
\begin{ex}
% ID: L12C2B13-13-DS
% Mức: VD
Xét dạng “Bài toán quá trình đẳng tích”: a) Bài mẫu cho kết quả 150. b) Kết quả bài mẫu gấp đôi 150. c) Lời giải phải nêu công thức trước khi thay số. d) Mọi dữ kiện số đều không cần đổi đơn vị.
\choiceTF

\loigiai{
a) Đúng: phù hợp với công thức và điều kiện đã nêu. b) Sai: phát biểu không thỏa công thức hoặc điều kiện áp dụng. c) Đúng: phù hợp với công thức và điều kiện đã nêu. d) Sai: phát biểu không thỏa công thức hoặc điều kiện áp dụng.
}
\end{ex}

% ===== Câu 38 =====
% ID: L12C2B13-13-TLN
% Mức: TH
\begin{ex}
% ID: L12C2B13-13-TLN
% Mức: TH
Khí trong bình kín cứng có áp suất $100\,\text{kPa}$ ở 300 K. Nung đến 450 K. Tính áp suất sau theo kPa.
\shortans{150}
\loigiai{
$p_2=p_1T_2/T_1=100\cdot450/300=150$ kPa.
}
\end{ex}

% ===== Câu 39 =====
% ID: L12C2B13-13-TN
% Mức: TH
\dangbt{Bài toán quá trình đẳng nhiệt}
\begin{ex}
% ID: L12C2B13-13-TN
% Mức: TH
Nhận xét đúng về quan hệ các đại lượng là
\choice
{các đại lượng không phụ thuộc nhau}
{mọi đại lượng đều không đổi}
{\True p tỉ lệ nghịch với $V$}
{chỉ phụ thuộc thời gian}
\loigiai{
Đáp án C. p tỉ lệ nghịch với V.
}
\end{ex}

% ===== Câu 40 =====
% ID: L12C2B13-14-DS
% Mức: VD
\dangbt{Bài toán quá trình đẳng tích}
\begin{ex}
% ID: L12C2B13-14-DS
% Mức: VD
Xét dạng “Bài toán quá trình đẳng tích”: a) Quan hệ cần dùng là p tỉ lệ thuận với $T$ tuyệt đối. b) Có thể suy ra quan hệ mà không xét điều kiện quá trình. c) Kết quả phải phù hợp chiều biến thiên vật lí. d) Áp suất dư luôn thay được áp suất tuyệt đối.
\choiceTF

\loigiai{
a) Đúng: phù hợp với công thức và điều kiện đã nêu. b) Sai: phát biểu không thỏa công thức hoặc điều kiện áp dụng. c) Đúng: phù hợp với công thức và điều kiện đã nêu. d) Sai: phát biểu không thỏa công thức hoặc điều kiện áp dụng.
}
\end{ex}

% ===== Câu 41 =====
% ID: L12C2B13-14-TLN
% Mức: TH
\begin{ex}
% ID: L12C2B13-14-TLN
% Mức: TH
Đổi 27 ° $C$ sang nhiệt độ tuyệt đối theo kelvin.
\shortans{300}
\loigiai{
$T=27+273=300$ K.
}
\end{ex}

% ===== Câu 42 =====
% ID: L12C2B13-14-TN
% Mức: TH
\dangbt{Bài toán quá trình đẳng nhiệt}
\begin{ex}
% ID: L12C2B13-14-TN
% Mức: TH
Sai lầm cần tránh khi giải dạng này là
\choice
{ghi đầy đủ dữ kiện}
{đổi đơn vị đúng}
{kiểm tra điều kiện}
{\True kết luận p tỉ lệ thuận với $V$}
\loigiai{
Đáp án D. Cần tránh: kết luận p tỉ lệ thuận với V.
}
\end{ex}

% ===== Câu 43 =====
% ID: L12C2B13-15-DS
% Mức: VDC
\dangbt{Bài toán quá trình đẳng tích}
\begin{ex}
% ID: L12C2B13-15-DS
% Mức: VDC
Xét dạng “Bài toán quá trình đẳng tích”: a) Dạng này có thể yêu cầu phối hợp đổi đơn vị và lập tỉ số. b) Kelvin được đổi sang Celsius bằng cách cộng 273. c) Cần ghi rõ đại lượng cần tìm. d) Công thức nào có đủ kí hiệu cũng dùng được.
\choiceTF

\loigiai{
a) Đúng: phù hợp với công thức và điều kiện đã nêu. b) Sai: phát biểu không thỏa công thức hoặc điều kiện áp dụng. c) Đúng: phù hợp với công thức và điều kiện đã nêu. d) Sai: phát biểu không thỏa công thức hoặc điều kiện áp dụng.
}
\end{ex}

% ===== Câu 44 =====
% ID: L12C2B13-15-TLN
% Mức: VD
\begin{ex}
% ID: L12C2B13-15-TLN
% Mức: VD
Một đại lượng tăng từ 100 lên 150. Tính tỉ số giá trị sau so với trước.
\shortans{1,5}
\loigiai{
Tỉ số bằng $150/100=1,5$.
}
\end{ex}

% ===== Câu 45 =====
% ID: L12C2B13-15-TN
% Mức: TH
\dangbt{Bài toán quá trình đẳng nhiệt}
\begin{ex}
% ID: L12C2B13-15-TN
% Mức: TH
Đơn vị thường dùng sau khi chuẩn hóa theo SI là
\choice
{\True Pa; m³}
{kg·m}
{W/s}
{C/mol}
\loigiai{
Đáp án A. Các đơn vị phù hợp: Pa; m³.
}
\end{ex}

% ===== Câu 46 =====
% ID: L12C2B13-16-DS
% Mức: TH
\dangbt{Bài toán quá trình đẳng áp}
\begin{ex}
% ID: L12C2B13-16-DS
% Mức: TH
Xét dạng “Bài toán quá trình đẳng áp”: a) Công thức đặc trưng là $V_1/T_1=V_2/T_2$. b) Có thể bỏ qua điều kiện lượng khí và áp suất không đổi. c) $V$ tỉ lệ thuận với $T$ tuyệt đối. d) Có thể dùng độ Celsius thay trực tiếp cho T.
\choiceTF

\loigiai{
a) Đúng: phù hợp với công thức và điều kiện đã nêu. b) Sai: phát biểu không thỏa công thức hoặc điều kiện áp dụng. c) Đúng: phù hợp với công thức và điều kiện đã nêu. d) Sai: phát biểu không thỏa công thức hoặc điều kiện áp dụng.
}
\end{ex}

% ===== Câu 47 =====
% ID: L12C2B13-16-TLN
% Mức: VD
\dangbt{Bài toán quá trình đẳng tích}
\begin{ex}
% ID: L12C2B13-16-TLN
% Mức: VD
Một thể tích 2,5 $L$ bằng bao nhiêu m³? Viết kết quả theo dạng $a\cdot10^{-3}$.
\shortans{2,5}
\loigiai{
$2,5L=2,5\cdot10^{-3}$ m³.
}
\end{ex}

% ===== Câu 48 =====
% ID: L12C2B13-16-TN
% Mức: VD
\dangbt{Bài toán quá trình đẳng nhiệt}
\begin{ex}
% ID: L12C2B13-16-TN
% Mức: VD
Bước đầu tiên hợp lí khi giải bài là
\choice
{thay số ngay}
{\True xác định quá trình và đại lượng không đổi}
{bỏ qua đơn vị}
{đổi áp suất thành nhiệt độ}
\loigiai{
Đáp án B. Trước hết phải nhận diện quá trình, liệt kê dữ kiện và đại lượng không đổi.
}
\end{ex}

% ===== Câu 49 =====
% ID: L12C2B13-17-DS
% Mức: TH
\dangbt{Bài toán quá trình đẳng áp}
\begin{ex}
% ID: L12C2B13-17-DS
% Mức: TH
Xét dạng “Bài toán quá trình đẳng áp”: a) Cần dùng đơn vị phù hợp: $K$; m³. b) Nhiệt độ tuyệt đối có thể mang giá trị âm. c) Phải xác định đại lượng không đổi trước khi tính. d) Không cần kiểm tra thứ nguyên.
\choiceTF

\loigiai{
a) Đúng: phù hợp với công thức và điều kiện đã nêu. b) Sai: phát biểu không thỏa công thức hoặc điều kiện áp dụng. c) Đúng: phù hợp với công thức và điều kiện đã nêu. d) Sai: phát biểu không thỏa công thức hoặc điều kiện áp dụng.
}
\end{ex}

% ===== Câu 50 =====
% ID: L12C2B13-17-TLN
% Mức: VD
\dangbt{Bài toán quá trình đẳng tích}
\begin{ex}
% ID: L12C2B13-17-TLN
% Mức: VD
Áp suất $120\,\text{kPa}$ bằng bao nhiêu Pa, lấy hệ số của $10^3$ làm đáp án?
\shortans{120}
\loigiai{
$120\,\text{kPa}=120\cdot10^3$ Pa.
}
\end{ex}

% ===== Câu 51 =====
% ID: L12C2B13-17-TN
% Mức: VD
\dangbt{Bài toán quá trình đẳng nhiệt}
\begin{ex}
% ID: L12C2B13-17-TN
% Mức: VD
Với dữ kiện của bài mẫu: Khí ở $100\,\text{kPa}$ có thể tích 6 $L$, bị nén đẳng nhiệt còn 4 L. Tính áp suất sau theo kPa.
\choice
{\True Kết quả 150}
{Kết quả 300}
{Không đủ dữ kiện}
{Kết quả bằng 0}
\loigiai{
Đáp án A. $p_2=p_1V_1/V_2=100\cdot6/4=150$ kPa.
}
\end{ex}

% ===== Câu 52 =====
% ID: L12C2B13-18-DS
% Mức: VD
\dangbt{Bài toán quá trình đẳng áp}
\begin{ex}
% ID: L12C2B13-18-DS
% Mức: VD
Xét dạng “Bài toán quá trình đẳng áp”: a) Bài mẫu cho kết quả 3. b) Kết quả bài mẫu gấp đôi 3. c) Lời giải phải nêu công thức trước khi thay số. d) Mọi dữ kiện số đều không cần đổi đơn vị.
\choiceTF

\loigiai{
a) Đúng: phù hợp với công thức và điều kiện đã nêu. b) Sai: phát biểu không thỏa công thức hoặc điều kiện áp dụng. c) Đúng: phù hợp với công thức và điều kiện đã nêu. d) Sai: phát biểu không thỏa công thức hoặc điều kiện áp dụng.
}
\end{ex}

% ===== Câu 53 =====
% ID: L12C2B13-18-TLN
% Mức: VDC
\dangbt{Bài toán quá trình đẳng tích}
\begin{ex}
% ID: L12C2B13-18-TLN
% Mức: VDC
Theo quan hệ “p tỉ lệ thuận với $T$ tuyệt đối”, nếu đại lượng phụ thuộc tỉ lệ thuận tăng gấp 4 thì hệ số tăng là bao nhiêu?
\shortans{4}
\loigiai{
Với quan hệ tỉ lệ thuận, hệ số biến đổi được giữ nguyên: 4.
}
\end{ex}

% ===== Câu 54 =====
% ID: L12C2B13-18-TN
% Mức: VD
\dangbt{Bài toán quá trình đẳng nhiệt}
\begin{ex}
% ID: L12C2B13-18-TN
% Mức: VD
Khi kiểm tra kết quả, tiêu chí vật lí quan trọng là
\choice
{chỉ cần số dương}
{không cần đơn vị}
{\True phù hợp chiều biến thiên theo định luật}
{luôn phải lớn hơn giá trị đầu}
\loigiai{
Đáp án C. Kết quả phải phù hợp với nhận xét: p tỉ lệ nghịch với V.
}
\end{ex}

% ===== Câu 55 =====
% ID: L12C2B13-19-DS
% Mức: VD
\dangbt{Bài toán quá trình đẳng áp}
\begin{ex}
% ID: L12C2B13-19-DS
% Mức: VD
Xét dạng “Bài toán quá trình đẳng áp”: a) Quan hệ cần dùng là $V$ tỉ lệ thuận với $T$ tuyệt đối. b) Có thể suy ra quan hệ mà không xét điều kiện quá trình. c) Kết quả phải phù hợp chiều biến thiên vật lí. d) Áp suất dư luôn thay được áp suất tuyệt đối.
\choiceTF

\loigiai{
a) Đúng: phù hợp với công thức và điều kiện đã nêu. b) Sai: phát biểu không thỏa công thức hoặc điều kiện áp dụng. c) Đúng: phù hợp với công thức và điều kiện đã nêu. d) Sai: phát biểu không thỏa công thức hoặc điều kiện áp dụng.
}
\end{ex}

% ===== Câu 56 =====
% ID: L12C2B13-19-TLN
% Mức: TH
\begin{ex}
% ID: L12C2B13-19-TLN
% Mức: TH
Khí có thể tích 2,0 $L$ ở 300 $K$, được nung đẳng áp đến 450 K. Tính thể tích sau theo L.
\shortans{3}
\loigiai{
$V_2=V_1T_2/T_1=2,0\cdot450/300=3,0$ L.
}
\end{ex}

% ===== Câu 57 =====
% ID: L12C2B13-19-TN
% Mức: VDC
\dangbt{Bài toán quá trình đẳng nhiệt}
\begin{ex}
% ID: L12C2B13-19-TN
% Mức: VDC
Nếu giữ đúng các điều kiện của định luật, biểu thức nên dùng là
\choice
{$s=vt$}
{\True $p_1V_1=p_2V_2$}
{$P=A/t$}
{$F=ma$}
\loigiai{
Đáp án B. Biểu thức đặc trưng là $p_1V_1=p_2V_2$.
}
\end{ex}

% ===== Câu 58 =====
% ID: L12C2B13-20-DS
% Mức: VDC
\dangbt{Bài toán quá trình đẳng áp}
\begin{ex}
% ID: L12C2B13-20-DS
% Mức: VDC
Xét dạng “Bài toán quá trình đẳng áp”: a) Dạng này có thể yêu cầu phối hợp đổi đơn vị và lập tỉ số. b) Kelvin được đổi sang Celsius bằng cách cộng 273. c) Cần ghi rõ đại lượng cần tìm. d) Công thức nào có đủ kí hiệu cũng dùng được.
\choiceTF

\loigiai{
a) Đúng: phù hợp với công thức và điều kiện đã nêu. b) Sai: phát biểu không thỏa công thức hoặc điều kiện áp dụng. c) Đúng: phù hợp với công thức và điều kiện đã nêu. d) Sai: phát biểu không thỏa công thức hoặc điều kiện áp dụng.
}
\end{ex}

% ===== Câu 59 =====
% ID: L12C2B13-20-TLN
% Mức: TH
\begin{ex}
% ID: L12C2B13-20-TLN
% Mức: TH
Đổi 27 ° $C$ sang nhiệt độ tuyệt đối theo kelvin.
\shortans{300}
\loigiai{
$T=27+273=300$ K.
}
\end{ex}

% ===== Câu 60 =====
% ID: L12C2B13-20-TN
% Mức: VDC
\dangbt{Bài toán quá trình đẳng nhiệt}
\begin{ex}
% ID: L12C2B13-20-TN
% Mức: VDC
Phát biểu nào phản ánh đúng bản chất bài toán?
\choice
{Có thể bỏ qua nhiệt độ tuyệt đối}
{Không cần xác định lượng khí}
{Mọi quá trình đều đẳng nhiệt}
{\True Phải dùng $p_1V_1=p_2V_2$ và kiểm tra lượng khí và nhiệt độ không đổi}
\loigiai{
Đáp án D. Phải dùng $p_1V_1=p_2V_2$, đồng thời kiểm tra lượng khí và nhiệt độ không đổi.
}
\end{ex}

% ===== Câu 61 =====
% ID: L12C2B13-21-DS
% Mức: TH
\dangbt{Bài toán tổng hợp, đồ thị và thực tiễn về khí lí tưởng}
\begin{ex}
% ID: L12C2B13-21-DS
% Mức: TH
Xét dạng “Bài toán tổng hợp, đồ thị và thực tiễn về khí lí tưởng”: a) Công thức đặc trưng là kết hợp $pV=nRT$ với các định luật Boyle, Charles và đẳng tích. b) Có thể bỏ qua điều kiện chia quá trình thành từng giai đoạn và dùng áp suất tuyệt đối. c) mỗi đoạn đồ thị ứng với một ràng buộc trạng thái riêng. d) Có thể áp dụng một định luật đơn cho toàn bộ quá trình nhiều giai đoạn.
\choiceTF

\loigiai{
a) Đúng: phù hợp với công thức và điều kiện đã nêu. b) Sai: phát biểu không thỏa công thức hoặc điều kiện áp dụng. c) Đúng: phù hợp với công thức và điều kiện đã nêu. d) Sai: phát biểu không thỏa công thức hoặc điều kiện áp dụng.
}
\end{ex}

% ===== Câu 62 =====
% ID: L12C2B13-21-TLN
% Mức: VD
\dangbt{Bài toán quá trình đẳng áp}
\begin{ex}
% ID: L12C2B13-21-TLN
% Mức: VD
Một đại lượng tăng từ 100 lên 150. Tính tỉ số giá trị sau so với trước.
\shortans{1,5}
\loigiai{
Tỉ số bằng $150/100=1,5$.
}
\end{ex}

% ===== Câu 63 =====
% ID: L12C2B13-21-TN
% Mức: NB
\dangbt{Bài toán quá trình đẳng tích}
\begin{ex}
% ID: L12C2B13-21-TN
% Mức: NB
Công thức cốt lõi phù hợp nhất cho dạng “Bài toán quá trình đẳng tích” là
\choice
{\True $p_1/T_1=p_2/T_2$}
{$Q=mc\Delta T$}
{$A=Fs$}
{$U=RI$}
\loigiai{
Đáp án A. Dùng $p_1/T_1=p_2/T_2$ trong điều kiện lượng khí và thể tích không đổi.
}
\end{ex}

% ===== Câu 64 =====
% ID: L12C2B13-22-DS
% Mức: TH
\dangbt{Bài toán tổng hợp, đồ thị và thực tiễn về khí lí tưởng}
\begin{ex}
% ID: L12C2B13-22-DS
% Mức: TH
Xét dạng “Bài toán tổng hợp, đồ thị và thực tiễn về khí lí tưởng”: a) Cần dùng đơn vị phù hợp: $K$; Pa; m³. b) Nhiệt độ tuyệt đối có thể mang giá trị âm. c) Phải xác định đại lượng không đổi trước khi tính. d) Không cần kiểm tra thứ nguyên.
\choiceTF

\loigiai{
a) Đúng: phù hợp với công thức và điều kiện đã nêu. b) Sai: phát biểu không thỏa công thức hoặc điều kiện áp dụng. c) Đúng: phù hợp với công thức và điều kiện đã nêu. d) Sai: phát biểu không thỏa công thức hoặc điều kiện áp dụng.
}
\end{ex}

% ===== Câu 65 =====
% ID: L12C2B13-22-TLN
% Mức: VD
\dangbt{Bài toán quá trình đẳng áp}
\begin{ex}
% ID: L12C2B13-22-TLN
% Mức: VD
Một thể tích 2,5 $L$ bằng bao nhiêu m³? Viết kết quả theo dạng $a\cdot10^{-3}$.
\shortans{2,5}
\loigiai{
$2,5L=2,5\cdot10^{-3}$ m³.
}
\end{ex}

% ===== Câu 66 =====
% ID: L12C2B13-22-TN
% Mức: NB
\dangbt{Bài toán quá trình đẳng tích}
\begin{ex}
% ID: L12C2B13-22-TN
% Mức: NB
Điều kiện quan trọng khi áp dụng dạng “Bài toán quá trình đẳng tích” là
\choice
{khối lượng vật luôn bằng $1\,\text{kg}$}
{\True lượng khí và thể tích không đổi}
{nhiệt độ luôn bằng 0 ° $C$}
{áp suất luôn bằng áp suất khí quyển}
\loigiai{
Đáp án B. Điều kiện áp dụng là lượng khí và thể tích không đổi.
}
\end{ex}

% ===== Câu 67 =====
% ID: L12C2B13-23-DS
% Mức: VD
\dangbt{Bài toán tổng hợp, đồ thị và thực tiễn về khí lí tưởng}
\begin{ex}
% ID: L12C2B13-23-DS
% Mức: VD
Xét dạng “Bài toán tổng hợp, đồ thị và thực tiễn về khí lí tưởng”: a) Bài mẫu cho kết quả 100. b) Kết quả bài mẫu gấp đôi 100. c) Lời giải phải nêu công thức trước khi thay số. d) Mọi dữ kiện số đều không cần đổi đơn vị.
\choiceTF

\loigiai{
a) Đúng: phù hợp với công thức và điều kiện đã nêu. b) Sai: phát biểu không thỏa công thức hoặc điều kiện áp dụng. c) Đúng: phù hợp với công thức và điều kiện đã nêu. d) Sai: phát biểu không thỏa công thức hoặc điều kiện áp dụng.
}
\end{ex}

% ===== Câu 68 =====
% ID: L12C2B13-23-TLN
% Mức: VD
\dangbt{Bài toán quá trình đẳng áp}
\begin{ex}
% ID: L12C2B13-23-TLN
% Mức: VD
Áp suất $120\,\text{kPa}$ bằng bao nhiêu Pa, lấy hệ số của $10^3$ làm đáp án?
\shortans{120}
\loigiai{
$120\,\text{kPa}=120\cdot10^3$ Pa.
}
\end{ex}

% ===== Câu 69 =====
% ID: L12C2B13-23-TN
% Mức: TH
\dangbt{Bài toán quá trình đẳng tích}
\begin{ex}
% ID: L12C2B13-23-TN
% Mức: TH
Nhận xét đúng về quan hệ các đại lượng là
\choice
{các đại lượng không phụ thuộc nhau}
{mọi đại lượng đều không đổi}
{\True p tỉ lệ thuận với $T$ tuyệt đối}
{chỉ phụ thuộc thời gian}
\loigiai{
Đáp án C. p tỉ lệ thuận với $T$ tuyệt đối.
}
\end{ex}

% ===== Câu 70 =====
% ID: L12C2B13-24-DS
% Mức: VD
\dangbt{Bài toán tổng hợp, đồ thị và thực tiễn về khí lí tưởng}
\begin{ex}
% ID: L12C2B13-24-DS
% Mức: VD
Xét dạng “Bài toán tổng hợp, đồ thị và thực tiễn về khí lí tưởng”: a) Quan hệ cần dùng là mỗi đoạn đồ thị ứng với một ràng buộc trạng thái riêng. b) Có thể suy ra quan hệ mà không xét điều kiện quá trình. c) Kết quả phải phù hợp chiều biến thiên vật lí. d) Áp suất dư luôn thay được áp suất tuyệt đối.
\choiceTF

\loigiai{
a) Đúng: phù hợp với công thức và điều kiện đã nêu. b) Sai: phát biểu không thỏa công thức hoặc điều kiện áp dụng. c) Đúng: phù hợp với công thức và điều kiện đã nêu. d) Sai: phát biểu không thỏa công thức hoặc điều kiện áp dụng.
}
\end{ex}

% ===== Câu 71 =====
% ID: L12C2B13-24-TLN
% Mức: VDC
\dangbt{Bài toán quá trình đẳng áp}
\begin{ex}
% ID: L12C2B13-24-TLN
% Mức: VDC
Theo quan hệ “ $V$ tỉ lệ thuận với $T$ tuyệt đối”, nếu đại lượng phụ thuộc tỉ lệ thuận tăng gấp 4 thì hệ số tăng là bao nhiêu?
\shortans{4}
\loigiai{
Với quan hệ tỉ lệ thuận, hệ số biến đổi được giữ nguyên: 4.
}
\end{ex}

% ===== Câu 72 =====
% ID: L12C2B13-24-TN
% Mức: TH
\dangbt{Bài toán quá trình đẳng tích}
\begin{ex}
% ID: L12C2B13-24-TN
% Mức: TH
Sai lầm cần tránh khi giải dạng này là
\choice
{ghi đầy đủ dữ kiện}
{đổi đơn vị đúng}
{kiểm tra điều kiện}
{\True bỏ qua việc đổi nhiệt độ sang kelvin}
\loigiai{
Đáp án D. Cần tránh: bỏ qua việc đổi nhiệt độ sang kelvin.
}
\end{ex}

% ===== Câu 73 =====
% ID: L12C2B13-25-DS
% Mức: VDC
\dangbt{Bài toán tổng hợp, đồ thị và thực tiễn về khí lí tưởng}
\begin{ex}
% ID: L12C2B13-25-DS
% Mức: VDC
Xét dạng “Bài toán tổng hợp, đồ thị và thực tiễn về khí lí tưởng”: a) Dạng này có thể yêu cầu phối hợp đổi đơn vị và lập tỉ số. b) Kelvin được đổi sang Celsius bằng cách cộng 273. c) Cần ghi rõ đại lượng cần tìm. d) Công thức nào có đủ kí hiệu cũng dùng được.
\choiceTF

\loigiai{
a) Đúng: phù hợp với công thức và điều kiện đã nêu. b) Sai: phát biểu không thỏa công thức hoặc điều kiện áp dụng. c) Đúng: phù hợp với công thức và điều kiện đã nêu. d) Sai: phát biểu không thỏa công thức hoặc điều kiện áp dụng.
}
\end{ex}

% ===== Câu 74 =====
% ID: L12C2B13-25-TLN
% Mức: TH
\begin{ex}
% ID: L12C2B13-25-TLN
% Mức: TH
Khí từ ($100\,\text{kPa}$; 2 $L$; 300 K) đến trạng thái có $V$ = 3 $L$ và $T$ = 450 K. Tính áp suất cuối theo kPa.
\shortans{100}
\loigiai{
$p_2=p_1V_1T_2/(V_2T_1)=100\cdot2\cdot450/(3\cdot300)=100$ kPa.
}
\end{ex}

% ===== Câu 75 =====
% ID: L12C2B13-25-TN
% Mức: TH
\dangbt{Bài toán quá trình đẳng tích}
\begin{ex}
% ID: L12C2B13-25-TN
% Mức: TH
Đơn vị thường dùng sau khi chuẩn hóa theo SI là
\choice
{\True $K$; Pa}
{kg·m}
{W/s}
{C/mol}
\loigiai{
Đáp án A. Các đơn vị phù hợp: $K$; Pa.
}
\end{ex}

% ===== Câu 76 =====
% ID: L12C2B13-26-DS
% Mức: TH
\dangbt{Nhận diện và lựa chọn định luật khí phù hợp}
\begin{ex}
% ID: L12C2B13-26-DS
% Mức: TH
Xét dạng “Nhận diện và lựa chọn định luật khí phù hợp”: a) Công thức đặc trưng là Boyle: $pV$ không đổi; Charles: $V/T$ không đổi; đẳng tích: $p/T$ không đổi. b) Có thể bỏ qua điều kiện lượng khí không đổi và xác định đại lượng trạng thái được giữ cố định. c) chọn định luật theo đại lượng không đổi. d) Có thể dùng nhiệt độ Celsius trong các tỉ số trạng thái.
\choiceTF

\loigiai{
a) Đúng: phù hợp với công thức và điều kiện đã nêu. b) Sai: phát biểu không thỏa công thức hoặc điều kiện áp dụng. c) Đúng: phù hợp với công thức và điều kiện đã nêu. d) Sai: phát biểu không thỏa công thức hoặc điều kiện áp dụng.
}
\end{ex}

% ===== Câu 77 =====
% ID: L12C2B13-26-TLN
% Mức: TH
\dangbt{Bài toán tổng hợp, đồ thị và thực tiễn về khí lí tưởng}
\begin{ex}
% ID: L12C2B13-26-TLN
% Mức: TH
Đổi 27 ° $C$ sang nhiệt độ tuyệt đối theo kelvin.
\shortans{300}
\loigiai{
$T=27+273=300$ K.
}
\end{ex}

% ===== Câu 78 =====
% ID: L12C2B13-26-TN
% Mức: VD
\dangbt{Bài toán quá trình đẳng tích}
\begin{ex}
% ID: L12C2B13-26-TN
% Mức: VD
Bước đầu tiên hợp lí khi giải bài là
\choice
{thay số ngay}
{\True xác định quá trình và đại lượng không đổi}
{bỏ qua đơn vị}
{đổi áp suất thành nhiệt độ}
\loigiai{
Đáp án B. Trước hết phải nhận diện quá trình, liệt kê dữ kiện và đại lượng không đổi.
}
\end{ex}

% ===== Câu 79 =====
% ID: L12C2B13-27-DS
% Mức: TH
\dangbt{Nhận diện và lựa chọn định luật khí phù hợp}
\begin{ex}
% ID: L12C2B13-27-DS
% Mức: TH
Xét dạng “Nhận diện và lựa chọn định luật khí phù hợp”: a) Cần dùng đơn vị phù hợp: $K$; Pa; m³. b) Nhiệt độ tuyệt đối có thể mang giá trị âm. c) Phải xác định đại lượng không đổi trước khi tính. d) Không cần kiểm tra thứ nguyên.
\choiceTF

\loigiai{
a) Đúng: phù hợp với công thức và điều kiện đã nêu. b) Sai: phát biểu không thỏa công thức hoặc điều kiện áp dụng. c) Đúng: phù hợp với công thức và điều kiện đã nêu. d) Sai: phát biểu không thỏa công thức hoặc điều kiện áp dụng.
}
\end{ex}

% ===== Câu 80 =====
% ID: L12C2B13-27-TLN
% Mức: VD
\dangbt{Bài toán tổng hợp, đồ thị và thực tiễn về khí lí tưởng}
\begin{ex}
% ID: L12C2B13-27-TLN
% Mức: VD
Một đại lượng tăng từ 100 lên 150. Tính tỉ số giá trị sau so với trước.
\shortans{1,5}
\loigiai{
Tỉ số bằng $150/100=1,5$.
}
\end{ex}

% ===== Câu 81 =====
% ID: L12C2B13-27-TN
% Mức: VD
\dangbt{Bài toán quá trình đẳng tích}
\begin{ex}
% ID: L12C2B13-27-TN
% Mức: VD
Với dữ kiện của bài mẫu: Khí trong bình kín cứng có áp suất $100\,\text{kPa}$ ở 300 K. Nung đến 450 K. Tính áp suất sau theo kPa.
\choice
{\True Kết quả 150}
{Kết quả 300}
{Không đủ dữ kiện}
{Kết quả bằng 0}
\loigiai{
Đáp án A. $p_2=p_1T_2/T_1=100\cdot450/300=150$ kPa.
}
\end{ex}

% ===== Câu 82 =====
% ID: L12C2B13-28-DS
% Mức: VD
\dangbt{Nhận diện và lựa chọn định luật khí phù hợp}
\begin{ex}
% ID: L12C2B13-28-DS
% Mức: VD
Xét dạng “Nhận diện và lựa chọn định luật khí phù hợp”: a) Bài mẫu cho kết quả 360. b) Kết quả bài mẫu gấp đôi 360. c) Lời giải phải nêu công thức trước khi thay số. d) Mọi dữ kiện số đều không cần đổi đơn vị.
\choiceTF

\loigiai{
a) Đúng: phù hợp với công thức và điều kiện đã nêu. b) Sai: phát biểu không thỏa công thức hoặc điều kiện áp dụng. c) Đúng: phù hợp với công thức và điều kiện đã nêu. d) Sai: phát biểu không thỏa công thức hoặc điều kiện áp dụng.
}
\end{ex}

% ===== Câu 83 =====
% ID: L12C2B13-28-TLN
% Mức: VD
\dangbt{Bài toán tổng hợp, đồ thị và thực tiễn về khí lí tưởng}
\begin{ex}
% ID: L12C2B13-28-TLN
% Mức: VD
Một thể tích 2,5 $L$ bằng bao nhiêu m³? Viết kết quả theo dạng $a\cdot10^{-3}$.
\shortans{2,5}
\loigiai{
$2,5L=2,5\cdot10^{-3}$ m³.
}
\end{ex}

% ===== Câu 84 =====
% ID: L12C2B13-28-TN
% Mức: VD
\dangbt{Bài toán quá trình đẳng tích}
\begin{ex}
% ID: L12C2B13-28-TN
% Mức: VD
Khi kiểm tra kết quả, tiêu chí vật lí quan trọng là
\choice
{chỉ cần số dương}
{không cần đơn vị}
{\True phù hợp chiều biến thiên theo định luật}
{luôn phải lớn hơn giá trị đầu}
\loigiai{
Đáp án C. Kết quả phải phù hợp với nhận xét: p tỉ lệ thuận với $T$ tuyệt đối.
}
\end{ex}

% ===== Câu 85 =====
% ID: L12C2B13-29-DS
% Mức: VD
\dangbt{Nhận diện và lựa chọn định luật khí phù hợp}
\begin{ex}
% ID: L12C2B13-29-DS
% Mức: VD
Xét dạng “Nhận diện và lựa chọn định luật khí phù hợp”: a) Quan hệ cần dùng là chọn định luật theo đại lượng không đổi. b) Có thể suy ra quan hệ mà không xét điều kiện quá trình. c) Kết quả phải phù hợp chiều biến thiên vật lí. d) Áp suất dư luôn thay được áp suất tuyệt đối.
\choiceTF

\loigiai{
a) Đúng: phù hợp với công thức và điều kiện đã nêu. b) Sai: phát biểu không thỏa công thức hoặc điều kiện áp dụng. c) Đúng: phù hợp với công thức và điều kiện đã nêu. d) Sai: phát biểu không thỏa công thức hoặc điều kiện áp dụng.
}
\end{ex}

% ===== Câu 86 =====
% ID: L12C2B13-29-TLN
% Mức: VD
\dangbt{Bài toán tổng hợp, đồ thị và thực tiễn về khí lí tưởng}
\begin{ex}
% ID: L12C2B13-29-TLN
% Mức: VD
Áp suất $120\,\text{kPa}$ bằng bao nhiêu Pa, lấy hệ số của $10^3$ làm đáp án?
\shortans{120}
\loigiai{
$120\,\text{kPa}=120\cdot10^3$ Pa.
}
\end{ex}

% ===== Câu 87 =====
% ID: L12C2B13-29-TN
% Mức: VDC
\dangbt{Bài toán quá trình đẳng tích}
\begin{ex}
% ID: L12C2B13-29-TN
% Mức: VDC
Nếu giữ đúng các điều kiện của định luật, biểu thức nên dùng là
\choice
{$s=vt$}
{\True $p_1/T_1=p_2/T_2$}
{$P=A/t$}
{$F=ma$}
\loigiai{
Đáp án B. Biểu thức đặc trưng là $p_1/T_1=p_2/T_2$.
}
\end{ex}

% ===== Câu 88 =====
% ID: L12C2B13-30-DS
% Mức: VDC
\dangbt{Nhận diện và lựa chọn định luật khí phù hợp}
\begin{ex}
% ID: L12C2B13-30-DS
% Mức: VDC
Xét dạng “Nhận diện và lựa chọn định luật khí phù hợp”: a) Dạng này có thể yêu cầu phối hợp đổi đơn vị và lập tỉ số. b) Kelvin được đổi sang Celsius bằng cách cộng 273. c) Cần ghi rõ đại lượng cần tìm. d) Công thức nào có đủ kí hiệu cũng dùng được.
\choiceTF

\loigiai{
a) Đúng: phù hợp với công thức và điều kiện đã nêu. b) Sai: phát biểu không thỏa công thức hoặc điều kiện áp dụng. c) Đúng: phù hợp với công thức và điều kiện đã nêu. d) Sai: phát biểu không thỏa công thức hoặc điều kiện áp dụng.
}
\end{ex}

% ===== Câu 89 =====
% ID: L12C2B13-30-TLN
% Mức: VDC
\dangbt{Bài toán tổng hợp, đồ thị và thực tiễn về khí lí tưởng}
\begin{ex}
% ID: L12C2B13-30-TLN
% Mức: VDC
Theo quan hệ “mỗi đoạn đồ thị ứng với một ràng buộc trạng thái riêng”, nếu đại lượng phụ thuộc tỉ lệ thuận tăng gấp 4 thì hệ số tăng là bao nhiêu?
\shortans{4}
\loigiai{
Với quan hệ tỉ lệ thuận, hệ số biến đổi được giữ nguyên: 4.
}
\end{ex}

% ===== Câu 90 =====
% ID: L12C2B13-30-TN
% Mức: VDC
\dangbt{Bài toán quá trình đẳng tích}
\begin{ex}
% ID: L12C2B13-30-TN
% Mức: VDC
Phát biểu nào phản ánh đúng bản chất bài toán?
\choice
{Có thể bỏ qua nhiệt độ tuyệt đối}
{Không cần xác định lượng khí}
{Mọi quá trình đều đẳng nhiệt}
{\True Phải dùng $p_1/T_1=p_2/T_2$ và kiểm tra lượng khí và thể tích không đổi}
\loigiai{
Đáp án D. Phải dùng $p_1/T_1=p_2/T_2$, đồng thời kiểm tra lượng khí và thể tích không đổi.
}
\end{ex}

% ===== Câu 91 =====
% ID: L12C2B13-31-TL
% Mức: TH
\dangbt{Nhận diện và lựa chọn định luật khí phù hợp}
\begin{ex}
% ID: L12C2B13-31-TL
% Mức: TH
Trình bày cơ sở lí thuyết và điều kiện áp dụng của dạng “Nhận diện và lựa chọn định luật khí phù hợp”.
\choiceTF

\loigiai{
Boyle: $pV$ không đổi; Charles: $V/T$ không đổi; đẳng tích: $p/T$ không đổi. Điều kiện: lượng khí không đổi và xác định đại lượng trạng thái được giữ cố định. Chuẩn hóa đơn vị $K$; Pa; m³; sau đó lập phương trình và kiểm tra kết quả.
}
\end{ex}

% ===== Câu 92 =====
% ID: L12C2B13-31-TLN
% Mức: TH
\begin{ex}
% ID: L12C2B13-31-TLN
% Mức: TH
Một lượng khí có thể tích không đổi, áp suất tăng từ $100\,\text{kPa}$ lên $120\,\text{kPa}$, nhiệt độ đầu 300 K. Tính nhiệt độ sau.
\shortans{360}
\loigiai{
Đẳng tích: $p_1/T_1=p_2/T_2$, nên $T_2=300\cdot120/100=360$ K.
}
\end{ex}

% ===== Câu 93 =====
% ID: L12C2B13-31-TN
% Mức: NB
\dangbt{Bài toán quá trình đẳng áp}
\begin{ex}
% ID: L12C2B13-31-TN
% Mức: NB
Công thức cốt lõi phù hợp nhất cho dạng “Bài toán quá trình đẳng áp” là
\choice
{\True $V_1/T_1=V_2/T_2$}
{$Q=mc\Delta T$}
{$A=Fs$}
{$U=RI$}
\loigiai{
Đáp án A. Dùng $V_1/T_1=V_2/T_2$ trong điều kiện lượng khí và áp suất không đổi.
}
\end{ex}

% ===== Câu 94 =====
% ID: L12C2B13-32-TL
% Mức: TH
\dangbt{Nhận diện và lựa chọn định luật khí phù hợp}
\begin{ex}
% ID: L12C2B13-32-TL
% Mức: TH
Nêu quy trình giải một bài thuộc dạng “Nhận diện và lựa chọn định luật khí phù hợp”.
\choiceTF

\loigiai{
Bước 1: tóm tắt và đổi đơn vị. Bước 2: nhận diện lượng khí không đổi và xác định đại lượng trạng thái được giữ cố định. Bước 3: dùng Boyle: $pV$ không đổi; Charles: $V/T$ không đổi; đẳng tích: $p/T$ không đổi. Bước 4: giải, ghi đơn vị và đối chiếu chọn định luật theo đại lượng không đổi.
}
\end{ex}

% ===== Câu 95 =====
% ID: L12C2B13-32-TLN
% Mức: TH
\begin{ex}
% ID: L12C2B13-32-TLN
% Mức: TH
Đổi 27 ° $C$ sang nhiệt độ tuyệt đối theo kelvin.
\shortans{300}
\loigiai{
$T=27+273=300$ K.
}
\end{ex}

% ===== Câu 96 =====
% ID: L12C2B13-32-TN
% Mức: NB
\dangbt{Bài toán quá trình đẳng áp}
\begin{ex}
% ID: L12C2B13-32-TN
% Mức: NB
Điều kiện quan trọng khi áp dụng dạng “Bài toán quá trình đẳng áp” là
\choice
{khối lượng vật luôn bằng $1\,\text{kg}$}
{\True lượng khí và áp suất không đổi}
{nhiệt độ luôn bằng 0 ° $C$}
{áp suất luôn bằng áp suất khí quyển}
\loigiai{
Đáp án B. Điều kiện áp dụng là lượng khí và áp suất không đổi.
}
\end{ex}

% ===== Câu 97 =====
% ID: L12C2B13-33-TL
% Mức: VD
\dangbt{Nhận diện và lựa chọn định luật khí phù hợp}
\begin{ex}
% ID: L12C2B13-33-TL
% Mức: VD
Một lượng khí có thể tích không đổi, áp suất tăng từ $100\,\text{kPa}$ lên $120\,\text{kPa}$, nhiệt độ đầu 300 K. Tính nhiệt độ sau.
\choiceTF

\loigiai{
Đẳng tích: $p_1/T_1=p_2/T_2$, nên $T_2=300\cdot120/100=360$ K.
}
\end{ex}

% ===== Câu 98 =====
% ID: L12C2B13-33-TLN
% Mức: VD
\begin{ex}
% ID: L12C2B13-33-TLN
% Mức: VD
Một đại lượng tăng từ 100 lên 150. Tính tỉ số giá trị sau so với trước.
\shortans{1,5}
\loigiai{
Tỉ số bằng $150/100=1,5$.
}
\end{ex}

% ===== Câu 99 =====
% ID: L12C2B13-33-TN
% Mức: TH
\dangbt{Bài toán quá trình đẳng áp}
\begin{ex}
% ID: L12C2B13-33-TN
% Mức: TH
Nhận xét đúng về quan hệ các đại lượng là
\choice
{các đại lượng không phụ thuộc nhau}
{mọi đại lượng đều không đổi}
{\True $V$ tỉ lệ thuận với $T$ tuyệt đối}
{chỉ phụ thuộc thời gian}
\loigiai{
Đáp án C. $V$ tỉ lệ thuận với $T$ tuyệt đối.
}
\end{ex}

% ===== Câu 100 =====
% ID: L12C2B13-34-TL
% Mức: VD
\dangbt{Nhận diện và lựa chọn định luật khí phù hợp}
\begin{ex}
% ID: L12C2B13-34-TL
% Mức: VD
Giải thích vì sao cần tránh sai lầm: “dùng nhiệt độ Celsius trong các tỉ số trạng thái”.
\choiceTF

\loigiai{
Sai lầm này làm dữ kiện không cùng hệ đơn vị hoặc không đúng bản chất đại lượng, khiến phép thế vào Boyle: $pV$ không đổi; Charles: $V/T$ không đổi; đẳng tích: $p/T$ không đổi cho kết quả sai. Cần chuẩn hóa trước khi tính.
}
\end{ex}

% ===== Câu 101 =====
% ID: L12C2B13-34-TLN
% Mức: VD
\begin{ex}
% ID: L12C2B13-34-TLN
% Mức: VD
Một thể tích 2,5 $L$ bằng bao nhiêu m³? Viết kết quả theo dạng $a\cdot10^{-3}$.
\shortans{2,5}
\loigiai{
$2,5L=2,5\cdot10^{-3}$ m³.
}
\end{ex}

% ===== Câu 102 =====
% ID: L12C2B13-34-TN
% Mức: TH
\dangbt{Bài toán quá trình đẳng áp}
\begin{ex}
% ID: L12C2B13-34-TN
% Mức: TH
Sai lầm cần tránh khi giải dạng này là
\choice
{ghi đầy đủ dữ kiện}
{đổi đơn vị đúng}
{kiểm tra điều kiện}
{\True dùng độ Celsius thay trực tiếp cho $T$}
\loigiai{
Đáp án D. Cần tránh: dùng độ Celsius thay trực tiếp cho T.
}
\end{ex}

% ===== Câu 103 =====
% ID: L12C2B13-35-TL
% Mức: VD
\dangbt{Nhận diện và lựa chọn định luật khí phù hợp}
\begin{ex}
% ID: L12C2B13-35-TL
% Mức: VD
Phân tích cách kiểm tra tính hợp lí của kết quả trong dạng “Nhận diện và lựa chọn định luật khí phù hợp”.
\choiceTF

\loigiai{
Kiểm tra thứ nguyên theo $K$; Pa; m³; kiểm tra dấu và bậc lớn; cuối cùng đối chiếu xu hướng chọn định luật theo đại lượng không đổi. Nếu trái xu hướng phải xem lại điều kiện và phép đổi đơn vị.
}
\end{ex}

% ===== Câu 104 =====
% ID: L12C2B13-35-TLN
% Mức: VD
\begin{ex}
% ID: L12C2B13-35-TLN
% Mức: VD
Áp suất $120\,\text{kPa}$ bằng bao nhiêu Pa, lấy hệ số của $10^3$ làm đáp án?
\shortans{120}
\loigiai{
$120\,\text{kPa}=120\cdot10^3$ Pa.
}
\end{ex}

% ===== Câu 105 =====
% ID: L12C2B13-35-TN
% Mức: TH
\dangbt{Bài toán quá trình đẳng áp}
\begin{ex}
% ID: L12C2B13-35-TN
% Mức: TH
Đơn vị thường dùng sau khi chuẩn hóa theo SI là
\choice
{\True $K$; m³}
{kg·m}
{W/s}
{C/mol}
\loigiai{
Đáp án A. Các đơn vị phù hợp: $K$; m³.
}
\end{ex}

% ===== Câu 106 =====
% ID: L12C2B13-36-TL
% Mức: VDC
\dangbt{Nhận diện và lựa chọn định luật khí phù hợp}
\begin{ex}
% ID: L12C2B13-36-TL
% Mức: VDC
Đề xuất cách giải khi bài toán có nhiều giai đoạn liên quan đến dạng “Nhận diện và lựa chọn định luật khí phù hợp”.
\choiceTF

\loigiai{
Chia quá trình thành các trạng thái liên tiếp, ghi p, $V$, $T$ cho từng trạng thái; áp dụng Boyle: $pV$ không đổi; Charles: $V/T$ không đổi; đẳng tích: $p/T$ không đổi ở đoạn thỏa lượng khí không đổi và xác định đại lượng trạng thái được giữ cố định; dùng kết quả đoạn trước làm dữ kiện đoạn sau và kiểm tra tính liên tục.
}
\end{ex}

% ===== Câu 107 =====
% ID: L12C2B13-36-TLN
% Mức: VDC
\begin{ex}
% ID: L12C2B13-36-TLN
% Mức: VDC
Theo quan hệ “chọn định luật theo đại lượng không đổi”, nếu đại lượng phụ thuộc tỉ lệ thuận tăng gấp 4 thì hệ số tăng là bao nhiêu?
\shortans{4}
\loigiai{
Với quan hệ tỉ lệ thuận, hệ số biến đổi được giữ nguyên: 4.
}
\end{ex}

% ===== Câu 108 =====
% ID: L12C2B13-36-TN
% Mức: VD
\dangbt{Bài toán quá trình đẳng áp}
\begin{ex}
% ID: L12C2B13-36-TN
% Mức: VD
Bước đầu tiên hợp lí khi giải bài là
\choice
{thay số ngay}
{\True xác định quá trình và đại lượng không đổi}
{bỏ qua đơn vị}
{đổi áp suất thành nhiệt độ}
\loigiai{
Đáp án B. Trước hết phải nhận diện quá trình, liệt kê dữ kiện và đại lượng không đổi.
}
\end{ex}

% ===== Câu 109 =====
% ID: L12C2B13-37-TN
% Mức: VD
\begin{ex}
% ID: L12C2B13-37-TN
% Mức: VD
Với dữ kiện của bài mẫu: Khí có thể tích 2,0 $L$ ở 300 $K$, được nung đẳng áp đến 450 K. Tính thể tích sau theo L.
\choice
{\True Kết quả 3}
{Kết quả 6}
{Không đủ dữ kiện}
{Kết quả bằng 0}
\loigiai{
Đáp án A. $V_2=V_1T_2/T_1=2,0\cdot450/300=3,0$ L.
}
\end{ex}

% ===== Câu 110 =====
% ID: L12C2B13-38-TN
% Mức: VD
\begin{ex}
% ID: L12C2B13-38-TN
% Mức: VD
Khi kiểm tra kết quả, tiêu chí vật lí quan trọng là
\choice
{chỉ cần số dương}
{không cần đơn vị}
{\True phù hợp chiều biến thiên theo định luật}
{luôn phải lớn hơn giá trị đầu}
\loigiai{
Đáp án C. Kết quả phải phù hợp với nhận xét: $V$ tỉ lệ thuận với $T$ tuyệt đối.
}
\end{ex}

% ===== Câu 111 =====
% ID: L12C2B13-39-TN
% Mức: VDC
\begin{ex}
% ID: L12C2B13-39-TN
% Mức: VDC
Nếu giữ đúng các điều kiện của định luật, biểu thức nên dùng là
\choice
{$s=vt$}
{\True $V_1/T_1=V_2/T_2$}
{$P=A/t$}
{$F=ma$}
\loigiai{
Đáp án B. Biểu thức đặc trưng là $V_1/T_1=V_2/T_2$.
}
\end{ex}

% ===== Câu 112 =====
% ID: L12C2B13-40-TN
% Mức: VDC
\begin{ex}
% ID: L12C2B13-40-TN
% Mức: VDC
Phát biểu nào phản ánh đúng bản chất bài toán?
\choice
{Có thể bỏ qua nhiệt độ tuyệt đối}
{Không cần xác định lượng khí}
{Mọi quá trình đều đẳng nhiệt}
{\True Phải dùng $V_1/T_1=V_2/T_2$ và kiểm tra lượng khí và áp suất không đổi}
\loigiai{
Đáp án D. Phải dùng $V_1/T_1=V_2/T_2$, đồng thời kiểm tra lượng khí và áp suất không đổi.
}
\end{ex}

% ===== Câu 113 =====
% ID: L12C2B13-41-TN
% Mức: NB
\dangbt{Bài toán tổng hợp, đồ thị và thực tiễn về khí lí tưởng}
\begin{ex}
% ID: L12C2B13-41-TN
% Mức: NB
Công thức cốt lõi phù hợp nhất cho dạng “Bài toán tổng hợp, đồ thị và thực tiễn về khí lí tưởng” là
\choice
{\True kết hợp $pV=nRT$ với các định luật Boyle, Charles và đẳng tích}
{$Q=mc\Delta T$}
{$A=Fs$}
{$U=RI$}
\loigiai{
Đáp án A. Dùng kết hợp $pV=nRT$ với các định luật Boyle, Charles và đẳng tích trong điều kiện chia quá trình thành từng giai đoạn và dùng áp suất tuyệt đối.
}
\end{ex}

% ===== Câu 114 =====
% ID: L12C2B13-42-TN
% Mức: NB
\begin{ex}
% ID: L12C2B13-42-TN
% Mức: NB
Điều kiện quan trọng khi áp dụng dạng “Bài toán tổng hợp, đồ thị và thực tiễn về khí lí tưởng” là
\choice
{khối lượng vật luôn bằng $1\,\text{kg}$}
{\True chia quá trình thành từng giai đoạn và dùng áp suất tuyệt đối}
{nhiệt độ luôn bằng 0 ° $C$}
{áp suất luôn bằng áp suất khí quyển}
\loigiai{
Đáp án B. Điều kiện áp dụng là chia quá trình thành từng giai đoạn và dùng áp suất tuyệt đối.
}
\end{ex}

% ===== Câu 115 =====
% ID: L12C2B13-43-TN
% Mức: TH
\begin{ex}
% ID: L12C2B13-43-TN
% Mức: TH
Nhận xét đúng về quan hệ các đại lượng là
\choice
{các đại lượng không phụ thuộc nhau}
{mọi đại lượng đều không đổi}
{\True mỗi đoạn đồ thị ứng với một ràng buộc trạng thái riêng}
{chỉ phụ thuộc thời gian}
\loigiai{
Đáp án C. mỗi đoạn đồ thị ứng với một ràng buộc trạng thái riêng.
}
\end{ex}

% ===== Câu 116 =====
% ID: L12C2B13-44-TN
% Mức: TH
\begin{ex}
% ID: L12C2B13-44-TN
% Mức: TH
Sai lầm cần tránh khi giải dạng này là
\choice
{ghi đầy đủ dữ kiện}
{đổi đơn vị đúng}
{kiểm tra điều kiện}
{\True áp dụng một định luật đơn cho toàn bộ quá trình nhiều giai đoạn}
\loigiai{
Đáp án D. Cần tránh: áp dụng một định luật đơn cho toàn bộ quá trình nhiều giai đoạn.
}
\end{ex}

% ===== Câu 117 =====
% ID: L12C2B13-45-TN
% Mức: TH
\begin{ex}
% ID: L12C2B13-45-TN
% Mức: TH
Đơn vị thường dùng sau khi chuẩn hóa theo SI là
\choice
{\True $K$; Pa; m³}
{kg·m}
{W/s}
{C/mol}
\loigiai{
Đáp án A. Các đơn vị phù hợp: $K$; Pa; m³.
}
\end{ex}

% ===== Câu 118 =====
% ID: L12C2B13-46-TN
% Mức: VD
\begin{ex}
% ID: L12C2B13-46-TN
% Mức: VD
Bước đầu tiên hợp lí khi giải bài là
\choice
{thay số ngay}
{\True xác định quá trình và đại lượng không đổi}
{bỏ qua đơn vị}
{đổi áp suất thành nhiệt độ}
\loigiai{
Đáp án B. Trước hết phải nhận diện quá trình, liệt kê dữ kiện và đại lượng không đổi.
}
\end{ex}

% ===== Câu 119 =====
% ID: L12C2B13-47-TN
% Mức: VD
\begin{ex}
% ID: L12C2B13-47-TN
% Mức: VD
Với dữ kiện của bài mẫu: Khí từ ($100\,\text{kPa}$; 2 $L$; 300 K) đến trạng thái có $V$ = 3 $L$ và $T$ = 450 K. Tính áp suất cuối theo kPa.
\choice
{\True Kết quả 100}
{Kết quả 200}
{Không đủ dữ kiện}
{Kết quả bằng 0}
\loigiai{
Đáp án A. $p_2=p_1V_1T_2/(V_2T_1)=100\cdot2\cdot450/(3\cdot300)=100$ kPa.
}
\end{ex}

% ===== Câu 120 =====
% ID: L12C2B13-48-TN
% Mức: VD
\begin{ex}
% ID: L12C2B13-48-TN
% Mức: VD
Khi kiểm tra kết quả, tiêu chí vật lí quan trọng là
\choice
{chỉ cần số dương}
{không cần đơn vị}
{\True phù hợp chiều biến thiên theo định luật}
{luôn phải lớn hơn giá trị đầu}
\loigiai{
Đáp án C. Kết quả phải phù hợp với nhận xét: mỗi đoạn đồ thị ứng với một ràng buộc trạng thái riêng.
}
\end{ex}

% ===== Câu 121 =====
% ID: L12C2B13-49-TN
% Mức: VDC
\begin{ex}
% ID: L12C2B13-49-TN
% Mức: VDC
Nếu giữ đúng các điều kiện của định luật, biểu thức nên dùng là
\choice
{$s=vt$}
{\True kết hợp $pV=nRT$ với các định luật Boyle, Charles và đẳng tích}
{$P=A/t$}
{$F=ma$}
\loigiai{
Đáp án B. Biểu thức đặc trưng là kết hợp $pV=nRT$ với các định luật Boyle, Charles và đẳng tích.
}
\end{ex}

% ===== Câu 122 =====
% ID: L12C2B13-50-TN
% Mức: VDC
\begin{ex}
% ID: L12C2B13-50-TN
% Mức: VDC
Phát biểu nào phản ánh đúng bản chất bài toán?
\choice
{Có thể bỏ qua nhiệt độ tuyệt đối}
{Không cần xác định lượng khí}
{Mọi quá trình đều đẳng nhiệt}
{\True Phải dùng kết hợp $pV=nRT$ với các định luật Boyle, Charles và đẳng tích và kiểm tra chia quá trình thành từng giai đoạn và dùng áp suất tuyệt đối}
\loigiai{
Đáp án D. Phải dùng kết hợp $pV=nRT$ với các định luật Boyle, Charles và đẳng tích, đồng thời kiểm tra chia quá trình thành từng giai đoạn và dùng áp suất tuyệt đối.
}
\end{ex}

% ===== Câu 123 =====
% ID: L12C2B13-51-TN
% Mức: NB
\dangbt{Nhận diện và lựa chọn định luật khí phù hợp}
\begin{ex}
% ID: L12C2B13-51-TN
% Mức: NB
Công thức cốt lõi phù hợp nhất cho dạng “Nhận diện và lựa chọn định luật khí phù hợp” là
\choice
{\True Boyle: $pV$ không đổi; Charles: $V/T$ không đổi; đẳng tích: $p/T$ không đổi}
{$Q=mc\Delta T$}
{$A=Fs$}
{$U=RI$}
\loigiai{
Đáp án A. Dùng Boyle: $pV$ không đổi; Charles: $V/T$ không đổi; đẳng tích: $p/T$ không đổi trong điều kiện lượng khí không đổi và xác định đại lượng trạng thái được giữ cố định.
}
\end{ex}

% ===== Câu 124 =====
% ID: L12C2B13-52-TN
% Mức: NB
\begin{ex}
% ID: L12C2B13-52-TN
% Mức: NB
Điều kiện quan trọng khi áp dụng dạng “Nhận diện và lựa chọn định luật khí phù hợp” là
\choice
{khối lượng vật luôn bằng $1\,\text{kg}$}
{\True lượng khí không đổi và xác định đại lượng trạng thái được giữ cố định}
{nhiệt độ luôn bằng 0 ° $C$}
{áp suất luôn bằng áp suất khí quyển}
\loigiai{
Đáp án B. Điều kiện áp dụng là lượng khí không đổi và xác định đại lượng trạng thái được giữ cố định.
}
\end{ex}

% ===== Câu 125 =====
% ID: L12C2B13-53-TN
% Mức: TH
\begin{ex}
% ID: L12C2B13-53-TN
% Mức: TH
Nhận xét đúng về quan hệ các đại lượng là
\choice
{các đại lượng không phụ thuộc nhau}
{mọi đại lượng đều không đổi}
{\True chọn định luật theo đại lượng không đổi}
{chỉ phụ thuộc thời gian}
\loigiai{
Đáp án C. chọn định luật theo đại lượng không đổi.
}
\end{ex}

% ===== Câu 126 =====
% ID: L12C2B13-54-TN
% Mức: TH
\begin{ex}
% ID: L12C2B13-54-TN
% Mức: TH
Sai lầm cần tránh khi giải dạng này là
\choice
{ghi đầy đủ dữ kiện}
{đổi đơn vị đúng}
{kiểm tra điều kiện}
{\True dùng nhiệt độ Celsius trong các tỉ số trạng thái}
\loigiai{
Đáp án D. Cần tránh: dùng nhiệt độ Celsius trong các tỉ số trạng thái.
}
\end{ex}

% ===== Câu 127 =====
% ID: L12C2B13-55-TN
% Mức: TH
\begin{ex}
% ID: L12C2B13-55-TN
% Mức: TH
Đơn vị thường dùng sau khi chuẩn hóa theo SI là
\choice
{\True $K$; Pa; m³}
{kg·m}
{W/s}
{C/mol}
\loigiai{
Đáp án A. Các đơn vị phù hợp: $K$; Pa; m³.
}
\end{ex}

% ===== Câu 128 =====
% ID: L12C2B13-56-TN
% Mức: VD
\begin{ex}
% ID: L12C2B13-56-TN
% Mức: VD
Bước đầu tiên hợp lí khi giải bài là
\choice
{thay số ngay}
{\True xác định quá trình và đại lượng không đổi}
{bỏ qua đơn vị}
{đổi áp suất thành nhiệt độ}
\loigiai{
Đáp án B. Trước hết phải nhận diện quá trình, liệt kê dữ kiện và đại lượng không đổi.
}
\end{ex}

% ===== Câu 129 =====
% ID: L12C2B13-57-TN
% Mức: VD
\begin{ex}
% ID: L12C2B13-57-TN
% Mức: VD
Với dữ kiện của bài mẫu: Một lượng khí có thể tích không đổi, áp suất tăng từ $100\,\text{kPa}$ lên $120\,\text{kPa}$, nhiệt độ đầu 300 K. Tính nhiệt độ sau.
\choice
{\True Kết quả 360}
{Kết quả 720}
{Không đủ dữ kiện}
{Kết quả bằng 0}
\loigiai{
Đáp án A. Đẳng tích: $p_1/T_1=p_2/T_2$, nên $T_2=300\cdot120/100=360$ K.
}
\end{ex}

% ===== Câu 130 =====
% ID: L12C2B13-58-TN
% Mức: VD
\begin{ex}
% ID: L12C2B13-58-TN
% Mức: VD
Khi kiểm tra kết quả, tiêu chí vật lí quan trọng là
\choice
{chỉ cần số dương}
{không cần đơn vị}
{\True phù hợp chiều biến thiên theo định luật}
{luôn phải lớn hơn giá trị đầu}
\loigiai{
Đáp án C. Kết quả phải phù hợp với nhận xét: chọn định luật theo đại lượng không đổi.
}
\end{ex}

% ===== Câu 131 =====
% ID: L12C2B13-59-TN
% Mức: VDC
\begin{ex}
% ID: L12C2B13-59-TN
% Mức: VDC
Nếu giữ đúng các điều kiện của định luật, biểu thức nên dùng là
\choice
{$s=vt$}
{\True Boyle: $pV$ không đổi; Charles: $V/T$ không đổi; đẳng tích: $p/T$ không đổi}
{$P=A/t$}
{$F=ma$}
\loigiai{
Đáp án B. Biểu thức đặc trưng là Boyle: $pV$ không đổi; Charles: $V/T$ không đổi; đẳng tích: $p/T$ không đổi.
}
\end{ex}

% ===== Câu 132 =====
% ID: L12C2B13-60-TN
% Mức: VDC
\begin{ex}
% ID: L12C2B13-60-TN
% Mức: VDC
Phát biểu nào phản ánh đúng bản chất bài toán?
\choice
{Có thể bỏ qua nhiệt độ tuyệt đối}
{Không cần xác định lượng khí}
{Mọi quá trình đều đẳng nhiệt}
{\True Phải dùng Boyle: $pV$ không đổi; Charles: $V/T$ không đổi; đẳng tích: $p/T$ không đổi và kiểm tra lượng khí không đổi và xác định đại lượng trạng thái được giữ cố định}
\loigiai{
Đáp án D. Phải dùng Boyle: $pV$ không đổi; Charles: $V/T$ không đổi; đẳng tích: $p/T$ không đổi, đồng thời kiểm tra lượng khí không đổi và xác định đại lượng trạng thái được giữ cố định.
}
\end{ex}
